% SPDX-License-Identifier: CC-BY-4.0
% Copyright (c) 2026 Martin Schröder <info@swedishembedded.com>

\section{Experimental design}\label{sec:design}

Four experiments exercise the flywheel. All are small, all were run once,
and none was tuned on its own measurement; they establish that the machinery behaves as
specified and give point estimates with their uncertainty, not a benchmark. Hardware
and software versions are in \cref{sec:repro}.

\paragraph{E1: corpus-expansion funnel on a real corpus.} The corpus is three Markdown pages of Brain's user documentation (monitoring, security and the D-Bus API; 10625 bytes) captured as one directory source. We ask a generator model (Qwen3-4B) for tasks
of four kinds, \code{recall}, \code{explain}, \code{classify} and the code-built
\code{denoise}. We report from the task report the number of proposals, admissions and
rejections per kind and per reason, and the coverage of the corpus's text parts, as the
funnel of \cref{sec:admission}.

\paragraph{E2: one full turn on a controlled domain.} The domain is 24 invented facts
per \emph{generation} (a bridge opened in a given year, a river of a given length, etc.;
every subject is an invented name and every value a random number, drawn from a seed
recorded with the run). A fact sheet is the source. The policy is Qwen3-0.6B, the
generator and teacher are Qwen3-4B. The turn is \code{splinter learn} with \code{recall}
tasks: tasks, closed-book rollouts of the policy (pass@$k$), a teacher solve where the
policy fails, verified experience, an \code{sft-final} dataset, LoRA training from the
champion, the four-check gate, release, alias move. The measurement is outside the
learner's reach: four probes per fact, worded by code in forms that no sheet or task
generator sees, asked of the live policy \code{policy:default} before the turn and after;
a fact counts as \emph{answered} when at least three of its four probes carry the right
value. A third set of facts that is never taught measures what the trained policy answers
about invented facts it was not given. We then roll the alias back and ask again.

\paragraph{E3: a second generation with retention.} The second fact sheet is to be taught by a
second \code{learn} that continues the champion with replay. We measure how many
first-generation facts remain answered (retention), the anchor suite of 40 general
questions, and the gate's own decision. Retention is the difference between the
first-generation answered count after the first release and after the second.

\paragraph{E4: the training substrate.} \Cref{sec:substrate} states what Brain's tree
contains. Whether the 27B model trains on one card in bf16-LoRA and NF4-QLoRA is a
measurement taken by Brain's own qualification gates; \cref{sec:e4} reports only what
their logs show.

\paragraph{Statistics.} Counts of answered facts carry Wilson 95\% intervals (the facts
are the unit; probes of one fact are not independent). Before-and-after comparisons on the
same facts use an exact two-sided McNemar test on discordant facts. With 24 facts per
sheet these are small samples; every E2 and E3 result is labelled a pilot.
